\section{An upper bound and two exact recurrences}\label{sec:upper}

We compare the source and target multiplicities without computing both
at every partition. This section supplies a gap-monotone upper bound
\(U(d)\ge F(d)\), together with exact recurrences for the target tables
and for the source values needed when that bound is inconclusive.
All three constructions come from coefficient extraction with Schur
alternants.

\subsection{The signed strip rule}
The coefficient rule below is the complete-symmetric case of the
plethystic Murnaghan--Nakayama rule; see
\cite[pp.~28--29]{DesarmenienLeclercThibon1994} and
\cite[arXiv:1408.3554v2, equation~(2)]{Wildon2016}.
Its residue-class determinant formulation also appears in
\cite[Theorem~3.4 and Corollary~3.11]{CaoJingLiu2025}.
We include the proof and the containment inequalities in the form
needed by the exact calculation.

For \(n\) variables put \(\delta_n=(n-1,\ldots,0)\).
The next rule bounds a single removal step in a product of factors
\(h_b(X^i)\), and also enumerates that step exactly.

\begin{lemma}[Signed strips]\label{lem:strips}
Let \(\gamma,\mu\) be partitions of length at most \(n\), and suppose
\(|\gamma|-|\mu|=ib\), where \(i,b\ge1\). Then
\[
                  [s_\gamma]\,h_b(X^i)s_\mu(X)\in\{0,1,-1\}.
\]
For a nonzero coefficient one necessarily has
\[
        \gamma_j\ge\mu_j,\qquad \mu_j\ge\gamma_{j+i}
                         \quad\text{whenever the subscripts are in range}.
\]
Here \(X^i=(x_1^i,\ldots,x_n^i)\).
\end{lemma}
\begin{proof}
Set \(l=\gamma+\delta_n\) and \(a=\mu+\delta_n\).
Multiplying by the Schur alternant denominator
\cite[Chapter~I, equation~(3.1)]{Macdonald1995} gives
\[
 [s_\gamma]h_b(X^i)s_\mu(X)
 =\det\bigl[\one_{\{a_k\le l_j,\ a_k\equiv l_j\pmod i\}}\bigr]_{j,k=1}^n.
\]
Indeed its determinant expansion assigns the entries of \(a\) to
positions \(j\), with assigned values \(\beta_j\) satisfying
\[
                  \beta_j\le l_j,\qquad \beta_j\equiv l_j\pmod i.
\]
The size condition makes the sum of the decrement quotients equal to
\(b\), so there is no further constraint from \(h_b(X^i)\).

Group the rows and columns by residue modulo \(i\), preserving decreasing
order within each group. Unequal row and column counts in a residue give
zero, and empty blocks contribute one. In a square block of size \(m\),
each row support is a terminal interval of columns, and the supports
decrease with the row index. A zero row or repeated support makes the
determinant zero. Otherwise their sizes are \(m,m-1,\ldots,1\): the block
is upper triangular with diagonal entries one, and its only allowed
matching is the diagonal. Translating this matching back to the original
positions gives
\begin{equation}\label{eq:stripconditions}
 \beta_j\le l_j,\qquad \beta_j>l_{j'}
      \quad\text{if \(j'\) is the next position of the same residue}.
\end{equation}
There is therefore just one surviving assignment. The sign from sorting
the assembled \(\beta\)'s is the remaining row-and-column ordering sign,
which proves the asserted coefficient values.

Sorting the componentwise inequalities \(\beta_j\le l_j\) gives
\(\mu_j\le\gamma_j\). For the other inequality, consider the first \(j+i\)
entries of \(l\). At most one per residue lacks a preceding entry of
that residue. Each of the remaining at least \(j\) entries has a
predecessor position whose assigned value, by
\eqref{eq:stripconditions}, is at least \(l_{j+i}+i\).
The \(j\)-th largest assigned value is therefore at least
\(l_{j+i}+i\), which after removing the staircases says
\(\mu_j\ge\gamma_{j+i}\).
\end{proof}

\subsection{A monotone upper certificate}
Averaging the permutations of six tensor factors projects onto
\(\Sym^6(\Sym^b V)\). Taking traces gives the cycle formula
\begin{equation}\label{eq:cycle}
        h_6[h_b](X)=\frac1{720}\sum_{\pi\in\SSS_6}
                   \prod_{\text{cycles }c\text{ of }\pi}h_b(X^{|c|}).
\end{equation}
We bound a coefficient of each product by removing its factors one at
a time. The signed strip rule restricts the intermediate partitions;
the bound below counts their possible coordinates while discarding
their signs and some of their constraints.

For a nonnegative gap vector \(d\), define
\[
 h_{sj}=1+\lambda(d)_j-\lambda(d)_{j+6-s},\qquad
 N_s(d)=\frac{\prod_{j=1}^s h_{sj}}{\max_{1\le j\le s}h_{sj}}
            \quad(2\le s\le5),\qquad N_0=N_1=1.
\]
For an ordered list \(I=(i_1,\ldots,i_v)\) with positive entries summing
to six put
\[
 B_I(d)=\prod_{u=1}^v N_{6-i_1-\cdots-i_u}(d),\qquad
 U(d)=\left\lfloor\frac1{720}
       \sum_{\pi\in\SSS_6}\min_{I\sim\pi}B_I(d)\right\rfloor ,
\]
where \(I\sim\pi\) means any ordering of the cycle lengths of \(\pi\).
This definition makes sense even when \(6\nmid w(d)\).

\begin{proposition}\label{prop:upper}
The function \(U\) is nondecreasing in every gap. For every nonnegative
gap vector \(d\) with \(6\mid w(d)\), one has \(F(d)\le U(d)\).
\end{proposition}
\begin{proof}
Each \(N_s\) is the minimum of the products obtained by omitting one of
the widths \(h_{sj}\). Every width is nondecreasing in the gaps, proving
monotonicity through the products, minima, sum, and floor.

The zero vector has \(F(0)=U(0)=1\), so assume \(b=w(d)/6\ge1\).
Remove the factors of one product in \eqref{eq:cycle} in an order \(I\).
After lengths totaling \(6-s\) have been removed, the remaining product
is the trace of a permutation of the \(s\) factors on
\((\Sym^b V)^{\otimes s}\), together with the diagonal \(\GL(V)\)-action.
The permutation commutes with that action, so its Schur support lies
among the constituents of the tensor power. Those have length at most
\(s\) by Pieri's rule
\cite[Chapter~I, equation~(5.16)]{Macdonald1995}.
Iterating Lemma~\ref{lem:strips}, every possible intermediate
partition \(\mu\) satisfies
\[
       |\mu|=sb,\qquad
       \lambda_{j+6-s}\le\mu_j\le\lambda_j\quad(1\le j\le s).
\]
Fix all but one coordinate; the sum fixes the remaining one.
Thus there are at most \(N_s(d)\) choices for \(2\le s\le5\), and at most
one for \(s=0,1\). Each chain has absolute coefficient at most one.
The coefficient of the product consequently has absolute value at most
\(B_I(d)\), for every order \(I\). Average these bounds in
\eqref{eq:cycle} and use the integrality of \(F(d)\).
\end{proof}

In Appendix~\ref{app:certificate}, \texttt{init\_masks} lists the eleven
cycle types with their class sizes and the remaining lengths
\(s\in\{2,3,4,5\}\) for each removal order. The routine
\texttt{ub\_cycle} forms the corresponding products of \(N_s\), takes
their minima, and performs the weighted sum and division by \(720\).

\subsection{Exact coefficients}
The target tables and the remaining source comparisons use Newton's
identity in place of a bound. For an integer vector \(v\) of length
\(n\), a \emph{signed lookup} is zero if \(v\) has a negative entry or a
repeated entry. Otherwise sort \(v\) decreasingly, subtract
\(\delta_n\), and multiply the coefficient at that partition by the
sorting sign. A strictly decreasing nonnegative integer list is at
least \(\delta_n\) componentwise, so this subtraction always gives a
partition.

Write \(Q_b(\lambda)=[s_\lambda]h_b[h_6]\) and
\(\widetilde Q_b(v)\) for its signed lookup in \(n\) variables.
The initial value is
\(Q_0(0)=1\). Newton's identity
\cite[Chapter~I, equation~(2.11)]{Macdonald1995}, multiplied by the alternating
denominator, gives
\begin{equation}\label{eq:targetrec}
 bQ_b(\lambda)=
 \sum_{i=1}^b\ \sum_{\substack{\alpha\in\NN^n\\|\alpha|=6}}
       \widetilde Q_{b-i}(\lambda+\delta_n-i\alpha).
\end{equation}
Indeed \(h_6(X^i)=\sum_{|\alpha|=6}X^{i\alpha}\).
Every nonzero previous partition \(\mu\) in this formula satisfies
\(\mu\le\lambda\) componentwise: subtracting a nonnegative vector cannot
increase any decreasing order statistic, and the same staircase is then
subtracted. In particular, a bound on a partition's tail is preserved in
all previous lookups.

For the source, fix \(b\) and set
\[
 H_j(\gamma)=[s_\gamma]h_j[h_b],\qquad H_0(0)=1.
\]
Newton's identity now gives the source recurrence; compare the
character formulation in
\cite[Proposition~5.1]{EvseevPagetWildon2014}:
\begin{equation}\label{eq:sourcerec}
 jH_j(\gamma)=\sum_{i=1}^j\ \sum_{\mu}
       [s_\gamma]\bigl(h_b(X^i)s_\mu(X)\bigr)H_{j-i}(\mu),
       \qquad |\mu|=(j-i)b.
\end{equation}
At stage \(j\), \(n=j\) variables suffice. Lemma~\ref{lem:strips} gives
an economical enumeration. Put \(l=\gamma+\delta_j\).
In each residue modulo \(i\), its last assigned
value must be that residue's least nonnegative representative: the
partition \(\mu\) has length at most \(j-i\), so its shifted list ends in
\(i-1,\ldots,0\). If any residue is absent, there is no term. At every
other position \(h\), whose next position in its residue is \(h'\), choose
\[
 \beta_h=l_h-ix_h,\qquad
               0\le x_h\le(l_h-l_{h'})/i-1.
\]
The sum of all decrement quotients, including the fixed last ones, is
\(b\). These choices have distinct values: different residues cannot
coincide, and in the same residue
\(\beta_h>l_{h'}\ge\beta_{h'}\). The fixed representatives
\(0,\ldots,i-1\) are the smallest \(i\) values. Thus sorting and
subtracting \(\delta_j\) gives a partition \(\mu\) of length at most
\(j-i\), with the sorting sign. Every such choice contributes once.
This is exactly the routine \texttt{strip} in Appendix~\ref{app:certificate}.
Equation~\eqref{eq:targetrec} is implemented by \texttt{gen},
\texttt{Make}, and \texttt{QTail}; \eqref{eq:sourcerec} by \texttt{Source}.
